\section{Conclusion}

ORION-Q demonstrates an executable way to preserve and recursively revisit negative computational research results without forcing every lane toward a positive. The method freezes outcome rules, gives donors first refusal, retains negative history, types the state that licenses reopening, and binds computations to replayable records. The claim-preserving recovery theorem adds a formal boundary: same-identity repair must causally cover every failed load-bearing predicate, while an uncovered predicate certifies that no repair exists within the declared action language. Identity-changing successes are successors rather than retrospective repairs.

Independent replay shows that the registered headline computations can be reproduced and that the replay harness can detect perturbation and missing generators. This strengthens the reproducibility claim, not the external-effectiveness claim.

That is the paper's claim ceiling. The work does not establish that this workflow is generally superior to alternative research methodologies, nor that a receipted result is automatically good evidence. Its value is a reproducible and formally bounded contract for saying what was tried, what failed, what a same-identity repair must change, why a question was reopened, and which stronger claim remains unauthorized.